\section*{Risk and Environmental Impact Assessment}
\setcounter{section}{6} \setcounter{subsection}{0}
\phantomsection \addcontentsline{toc}{section}{\textit{Risk and Environmental Impact Assessment}}
% Please use only \subsection*{} and \subsubsection*{} in this section. This suppresses the subsection numbering in the Table of Contents.


\textbf{Risk and Environmental Impact Assessment}

This project investigates \texttt{linux-mcp}, a kernel-assisted control-plane architecture for AI tool invocation.
Its goal is to improve the safety and observability of tool use in AI-agent systems by separating semantic interpretation and tool execution in user space from admission control and approval handling in the kernel.
Because the system spans an LLM-facing application, a user-space gateway, kernel mechanisms, and external tools, it introduces technical, security, and resource-related risks.
A formal risk and environmental impact assessment is therefore necessary.

Following the prescribed method, each risk is assigned a likelihood level $L$ and a consequence level $C$, and the overall score is calculated as $R = L \cdot C$.
This assessment considers risks to project completion, risks of harm to systems or organisations, environmental impact, and losses of time or computational resources.

\textbf{Project and technical risks}

A primary technical risk is kernel-level instability.
Since \texttt{linux-mcp} moves part of the control plane into the Linux kernel, faults in Netlink communication, synchronization, or state management may lead to crashes or inconsistent behaviour.
This risk is assigned moderate likelihood ($L = 3$) and major consequence ($C = 4$), giving $R = 12$.
The mitigation is to keep the kernel scope narrow and rely on strict interfaces, staged testing, and stress testing.

A second risk is the possibility that the project does not demonstrate a convincing advantage over strong user-space baselines.
The value of \texttt{linux-mcp} depends not only on correctness but also on showing that kernel placement provides stronger control properties than user-space mediation alone.
This risk is assigned likely likelihood ($L = 4$) and very serious consequence ($C = 3$), giving $R = 12$.
It should be mitigated by evaluation against both conventional and hardened user-space baselines under representative performance and adversarial workloads.

A third risk is cross-component integration failure.
The system includes an LLM-side application, the \texttt{mcpd} daemon, the kernel control plane, and tool-side services, and mismatches between these components may create failures that are difficult to isolate.
This risk is assigned moderate likelihood ($L = 3$) and serious consequence ($C = 2$), giving $R = 6$.
This can be reduced through explicit protocol definitions, schema synchronization, and end-to-end validation.

\textbf{Security and operational risks}

A central operational risk is unsafe or unauthorized tool invocation.
Since \texttt{linux-mcp} mediates LLM-triggered access to tools with possible side effects, an error in the admission path could allow an action that should have been denied or deferred.
This risk is assigned moderate likelihood ($L = 3$) and major consequence ($C = 4$), giving $R = 12$.
Mitigation includes deny-by-default admission, explicit approval binding, auditable state transitions, and attack-oriented testing.

A second risk is overestimating the system's security guarantees.
Kernel placement can strengthen enforcement of specific invariants, but it does not eliminate threats arising from semantic errors in user space, unsafe tool implementations, or incomplete policies.
This risk is assigned moderate likelihood ($L = 3$) and very serious consequence ($C = 3$), giving $R = 9$.
The mitigation is to define the threat model precisely and to state clearly which protections are provided by kernel enforcement and which remain outside its scope.

A further risk is degraded responsiveness under load.
Additional arbitration and approval steps may increase latency or reduce throughput, especially when multiple agents or tools are active concurrently.
This risk is assigned moderate likelihood ($L = 3$) and serious consequence ($C = 2$), giving $R = 6$.
It should be mitigated through latency analysis, concurrency benchmarking, and optimization of common low-risk paths.

\textbf{Environmental and resource impact}

The direct environmental impact of \texttt{linux-mcp} is limited because the project is software-based and does not involve substantial material waste.
However, development and evaluation still consume computational resources.
Repeated benchmarking, adversarial testing, virtual-machine execution, and debugging of kernel-space behaviour all increase CPU time and electricity consumption.
This environmental risk is assigned likely likelihood ($L = 4$) and serious consequence ($C = 2$), giving $R = 8$.
The appropriate mitigation is careful experiment planning, automation of result collection, and avoidance of redundant reruns.

A practical project-management risk is loss of development time and compute resources due to unstable experiments, failed integration, or repeated debugging.
This risk is assigned likely likelihood ($L = 4$) and serious consequence ($C = 2$), giving $R = 8$.
Mitigation includes reproducible scripts, checkpointing of experiment outputs, staged deployment, and careful configuration management.

Another resource-related risk is dependence on external tooling or infrastructure.
The project may rely on LLM APIs, network configuration, proxies, virtual machines, or other software dependencies during development and evaluation.
This risk is assigned moderate likelihood ($L = 3$) and serious consequence ($C = 2$), giving $R = 6$.
This risk can be reduced by separating the core evaluation from optional cloud dependencies and by supporting local or mocked execution paths where possible.

\begin{table}[t]
\centering
\caption{Summary of identified risks and mitigation strategies for \texttt{linux-mcp}.}
\begin{tabular}{p{4.1cm} c c c p{2.1cm} p{5.9cm}}
\hline
\textbf{Risk} & \textbf{$L$} & \textbf{$C$} & \textbf{$R$} & \textbf{Rating} & \textbf{Mitigation} \\
\hline
Kernel-level instability & 3 & 4 & 12 & Significant & Narrow kernel scope, strict interfaces, staged testing, stress testing \\
Weak justification over user-space baselines & 4 & 3 & 12 & Significant & Strong baselines, adversarial evaluation, quantitative comparison \\
Cross-component integration failure & 3 & 2 & 6 & Moderate & Explicit protocols, schema synchronization, end-to-end validation \\
Unsafe or unauthorized tool invocation & 3 & 4 & 12 & Significant & Deny-by-default admission, approval binding, auditability, misuse testing \\
Overestimated security guarantees & 3 & 3 & 9 & Significant & Precise threat model, careful claims, clear scope of guarantees \\
Added latency or reduced throughput & 3 & 2 & 6 & Moderate & Latency analysis, concurrency evaluation, fast-path optimization \\
Additional compute and energy use & 4 & 2 & 8 & Significant & Efficient experiment planning, automation, fewer redundant reruns \\
Loss of development time and compute resources & 4 & 2 & 8 & Significant & Reproducible workflows, checkpointing, staged deployment \\
External dependency or infrastructure failure & 3 & 2 & 6 & Moderate & Local or mocked backends, dependency isolation, reduced cloud reliance \\
\hline
\end{tabular}
\label{tab:risk_assessment}
\end{table}